\section{A evolução no tempo e a marca da pandemia}
\label{sec:tendencia}

As Figuras~\ref{fig:tendA} a~\ref{fig:tendC} acompanham a mediana anual
de cada indicador de 2015 a 2025, com a faixa cinza destacando o biênio
2020 e 2021. A pandemia deixou marca em quase tudo: a mortalidade (já sem as
internações COVID) sobe de cerca de 5,3\% para \textbf{5,85\% em 2020 e
6,28\% em 2021} e só retorna ao patamar histórico em 2025, sinal de que
o biênio afetou também os pacientes não COVID. A ocupação de internação
despenca para \textbf{64,5\% em 2020} (leitos esvaziados por
cancelamentos e medo) enquanto a de UTI dispara para \textbf{100,8\% em
2021}, seu recorde na série: a mediana da rede passa acima da
capacidade cadastrada, e não apenas a cauda, o que é a assinatura
aritmética da superlotação discutida na Seção~\ref{sec:qualidade}. O tempo de permanência atravessa tudo
estável, em torno de 4,3 dias.

O faturamento pede uma leitura em duas camadas. Em valores correntes
(nominais), a mediana sobe de R\$~866 em 2015 para R\$~1.283 em 2025.
Mas parte dessa alta é só inflação. Corrigindo tudo para o poder de
compra de 2025 com o IPCA (Figura~\ref{fig:custoreal} e
Tabela~\ref{tab:custoreal}), a história muda de sinal: \textbf{em
valores de 2025, o faturamento mediano caiu de R\$~1.427 para R\$~1.283, uma
queda real de cerca de 10\% na década}. Toda a alta aparente é
inflação. Qualquer avaliação da evolução do faturamento deve, portanto, ser
feita em valores reais ou com o nível de preços de cada ano devidamente
considerado.

Por categoria (Figuras~\ref{fig:tendcatA} a~\ref{fig:tendcatC}), o
choque pandêmico é
generalizado, mas a queda de ocupação de 2020 é mais funda na Direta e
nos Filantrópicos; a fração de alta complexidade das OSS cresce de
forma consistente na segunda metade do período; e, no faturamento em
valores correntes, a Direta
parte de nível parecido com as demais categorias e passa a liderar o
faturamento mediano a partir de 2020.

\begin{table}[H]
\centering
\caption{Faturamento por saída mediano por ano: em valores correntes e em
valores de 2025 (corrigidos pelo IPCA, variação de dezembro a dezembro,
IBGE). O fator indica por quanto o valor do ano é multiplicado para
chegar a preços de 2025.}
\label{tab:custoreal}
\small
\begin{tabular}{lrrr}
\toprule
Ano & Nominal (R\$) & Valores de 2025 (R\$) & Fator IPCA \\
\midrule
2015 & 866,15   & 1.427,31 & 1,6479 \\
2016 & 879,45   & 1.363,47 & 1,5504 \\
2017 & 882,36   & 1.328,77 & 1,5059 \\
2018 & 903,01   & 1.310,72 & 1,4515 \\
2019 & 934,26   & 1.300,04 & 1,3915 \\
2020 & 1.023,02 & 1.361,99 & 1,3313 \\
2021 & 1.050,43 & 1.270,66 & 1,2097 \\
2022 & 1.176,37 & 1.345,12 & 1,1435 \\
2023 & 1.197,21 & 1.308,50 & 1,0930 \\
2024 & 1.266,82 & 1.320,79 & 1,0426 \\
2025 & 1.283,46 & 1.283,46 & 1,0000 \\
\bottomrule
\end{tabular}
\end{table}

\comofeito{cada linha é um ano. A coluna \textbf{Nominal} é o faturamento
mediano por saída em reais correntes daquele ano; a coluna \textbf{Valores
de 2025} é o mesmo valor multiplicado pelo \textbf{Fator IPCA} (a
variação acumulada do IPCA de dezembro a dezembro, IBGE, do ano até
2025), trazendo tudo para o poder de compra de 2025. Em 2025 o fator é
1,0000 por definição, então as duas colunas coincidem.}
\newpage
\comoler{a linha é o hospital mediano em cada ano; a área azul em volta
contém a metade central dos hospitais (do quartil inferior ao
superior). Dentro da faixa cinza da pandemia, note a corcova da
mortalidade; o TMP, em contraste, atravessa a década quase imóvel.}

\begin{figure}[H]
\centering
\subfig{fig_ae_07_tendencia_mort_all.png}{Mortalidade geral}
\subfig{fig_ae_07_tendencia_tmp.png}{TMP}
\caption{Mediana anual da mortalidade e do TMP, com a metade central
dos hospitais na área sombreada em azul; a faixa cinza destaca 2020 e
2021 e os pontos de máximo, mínimo e último ano estão anotados.}
\label{fig:tendA}
\end{figure}

\comofeito{série temporal. O eixo horizontal é o ano (2015 a 2025) e o
vertical, o valor do indicador (mortalidade em proporção; TMP em dias). A
linha liga a \textbf{mediana} de todos os hospitais em cada ano; a área
sombreada em azul ao redor é a metade central dos hospitais (do quartil
inferior ao superior) e a faixa cinza destaca o biênio pandêmico de 2020
e 2021. Pontos de máximo, mínimo e último ano estão anotados.}
\newpage
\comoler{o faturamento real recua ao longo da década: descontada a
inflação, o valor pago por internação encerrada perde poder de compra. O
contraste direto com a linha nominal está na
Figura~\ref{fig:custoreal}. A alta complexidade mediana cresce nos anos
finais da série.}

\begin{figure}[H]
\centering
\subfig{fig_ae_07_tendencia_custo_real.png}{Faturamento real por saída
(R\$ de dez/2025)}
\subfig{fig_ae_07_tendencia_pct_alta_complex.png}{Fração alta
complexidade}
\caption{Mediana anual do faturamento real por saída (R\$ de dez/2025) e
da fração de alta complexidade.}
\label{fig:tendB}
\end{figure}

\comofeito{mesma construção da figura anterior (mediana anual e metade
central sombreada), para o faturamento real por saída (deflacionado pelo
IPCA para R\$ de dez/2025) e a fração de alta complexidade; ano no eixo
horizontal, indicador no vertical.}
\newpage
\comoler{dentro da faixa cinza da pandemia, os dois movimentos em
sentidos opostos: a internação esvazia em 2020 (mergulho para 64,5\%) e
a UTI lota em 2021 (pico de 86,1\%, recorde da série).}

\begin{figure}[H]
\centering
\subfig{fig_ae_07_tendencia_ocupacao_internacao.png}{Ocupação de
internação}
\subfig{fig_ae_07_tendencia_ocupacao_uti.png}{Ocupação de UTI}
\caption{Mediana anual das taxas de ocupação.}
\label{fig:tendC}
\end{figure}

\comofeito{mesma construção, para as duas taxas de ocupação (percentual)
no eixo vertical ao longo dos anos no eixo horizontal.}
\newpage
\comoler{as duas linhas contam histórias opostas. A clara (valores
correntes) sobe o tempo todo; a escura (valores de 2025) mostra o que o
dinheiro de fato compra e cai ao longo da década: de R\$~1.427 para
R\$~1.283 por internação. Quando as duas linhas se encontram, em 2025,
é porque ali o valor corrente e o corrigido coincidem por definição.}

\begin{figure}[H]
\centering
\includegraphics[width=.62\textwidth]{fig_ae_07_tendencia_custo_saida_real.png}
\caption{Faturamento por saída mediano: a linha clara tracejada é o valor
corrente e a linha escura é o valor corrigido pelo IPCA para o poder de
compra de 2025.}
\label{fig:custoreal}
\end{figure}

\comofeito{o eixo horizontal é o ano e o vertical, o faturamento mediano
por saída em reais. São desenhadas \textbf{duas linhas}: a clara
tracejada é o valor corrente (nominal) e a escura é o mesmo valor
corrigido pelo IPCA (variação de dezembro a dezembro, IBGE) para o poder
de compra de 2025. As duas coincidem em 2025 por definição, pois é o ano
de referência da correção.}
\newpage
\comoler{cada linha é o hospital mediano de uma categoria em cada ano.
Na mortalidade, todas as categorias fazem a corcova da pandemia. No
TMP, a linha da Direta corre acima das demais durante a década
inteira.}

\begin{figure}[H]
\centering
\subfig{fig_ae_08_categoria_mort_all.png}{Mortalidade geral}
\subfig{fig_ae_08_categoria_tmp.png}{TMP}
\caption{Mediana anual da mortalidade e do TMP por categoria de gestão;
cada categoria tem cor e marcador próprios, e a linha pontilhada cinza
do Privado (3 hospitais) é apenas registro.}
\label{fig:tendcatA}
\end{figure}

\comofeito{série temporal com uma linha por categoria de gestão. O eixo
horizontal é o ano e o vertical, a \textbf{mediana anual} do indicador
(mortalidade em proporção; TMP em dias) dentro de cada categoria. Cada
categoria tem cor e marcador próprios; a linha pontilhada cinza do
Privado (só 3 hospitais) aparece apenas como registro.}
\newpage
\comoler{no faturamento, as quatro categorias sobem juntas até 2019 e se
abrem a partir de 2020, com a Direta na frente. Na alta complexidade, a
linha das OSS se descola das demais na segunda metade do período.}

\begin{figure}[H]
\centering
\subfig{fig_ae_08_categoria_custo_real.png}{Faturamento real por saída,
R\$ de dez/2025 (Privado omitido por escala)}
\subfig{fig_ae_08_categoria_pct_alta_complex.png}{Fração alta
complexidade (Privado omitido por escala)}
\caption{Mediana anual do faturamento real por saída (R\$ de dez/2025) e
da fração de alta
complexidade por categoria. O Privado foi omitido destes dois painéis
porque sua escala (faturamento dez vezes maior) esconderia as demais
categorias; seus valores estão na Tabela~\ref{tab:categorias}.}
\label{fig:tendcatB}
\end{figure}

\comofeito{mesma construção da figura anterior (uma linha de mediana
anual por categoria), para o faturamento real por saída (R\$ de
dez/2025) e a fração de alta complexidade. O Privado foi omitido destes
dois painéis porque sua escala, dez vezes maior, esconderia as demais.}
\newpage
\comoler{na ocupação de internação, as duas linhas de cima (OSS e
Municipal) e as duas de baixo (Filantrópico e Direta) mantêm a
separação de quase 25 pontos de ocupação durante toda a década (as de
cima na casa dos 90\% e as de baixo na dos 65\%), com o
mergulho de 2020 visível em todas. Na UTI, o pico coletivo de 2021 é a
marca da pandemia.}

\begin{figure}[H]
\centering
\subfig{fig_ae_08_categoria_ocupacao_internacao.png}{Ocupação de
internação}
\subfig{fig_ae_08_categoria_ocupacao_uti.png}{Ocupação de UTI}
\caption{Mediana anual das taxas de ocupação por categoria de gestão.}
\label{fig:tendcatC}
\end{figure}

\comofeito{mesma construção, com uma linha de mediana anual por categoria
para as duas taxas de ocupação (percentual) no eixo vertical.}

